
\documentclass{article}
\usepackage{colortbl}
\usepackage{makecell}
\usepackage{multirow}
\usepackage{supertabular}

\begin{document}

\newcounter{utterance}

\centering \large Interaction Transcript for game `codenames', experiment `abstract', episode 8 with reflexion{-}llama{-}3.1{-}8b{-}instruct.
\vspace{24pt}

{ \footnotesize  \setcounter{utterance}{1}
\setlength{\tabcolsep}{0pt}
\begin{supertabular}{c@{$\;$}|p{.15\linewidth}@{}p{.15\linewidth}p{.15\linewidth}p{.15\linewidth}p{.15\linewidth}p{.15\linewidth}}
    \# & $\;$A & \multicolumn{4}{c}{Game Master} & $\;\:$B\\
    \hline

    \theutterance \stepcounter{utterance}  
    & & & \multicolumn{2}{p{0.3\linewidth}}{
        \cellcolor[rgb]{0.9,0.9,0.9}{
            \makecell[{{p{\linewidth}}}]{
                \texttt{\tiny{[GM$|$GM]}}
                {'hidden': {'team': ['dilution', 'reinstatement', 'gracefulness', 'nationalization'], 'innocent': ['stigma', 'ridiculousness', 'optimism', 'inefficiency', 'unsettlement', 'adorability'], 'opponent': [], 'assassin': ['coercion']}, 'revealed': {'team': {'team': [], 'innocent': [], 'opponent': [], 'assassin': []}, 'opponent': {'team': [], 'innocent': [], 'opponent': [], 'assassin': []}}}
            }
        }
    }
    & & \\ \\

    \theutterance \stepcounter{utterance}  
    & & \multicolumn{4}{p{0.6\linewidth}}{
        \cellcolor[rgb]{0.9,0.9,0.9}{
            \makecell[{{p{\linewidth}}}]{
                \texttt{\tiny{[P1$\langle$GM]}}
                \texttt{You have attempted this same task before and did not succeed. These are the plans you wrote after those attempts.} \\
\texttt{Trial \#1: New plan: } \\
\\ 
\texttt{In my previous attempts, I provided a single{-}word clue that was semantically related to some of my team's words. However, I failed to consider the potential for my clue to be related to an opponent's word, which led to my teammate's turn ending prematurely. } \\
\\ 
\texttt{Next time, I will make sure to carefully review the lists of opponent words and distractor words before providing a clue. I will also aim to create a clue that is more specific and less likely to be related to an opponent's word. For example, I could try to find a word that is related to a specific subset of my team's words, rather than trying to find a word that is related to multiple words. } \\
\\ 
\texttt{Additionally, I will pay closer attention to the format of my response, making sure to start with 'CLUE: ', followed by a new line, and then 'TARGETS: ' followed by a comma{-}separated list of related target words.} \\
\\ 
\texttt{Trial \#2: New plan: } \\
\\ 
\texttt{In my previous attempts, I provided a single{-}word clue that was semantically related to some of my team's words, but I failed to consider the potential for my clue to be related to an opponent's word, which led to my teammate's turn ending prematurely. } \\
\\ 
\texttt{Next time, I will make sure to carefully review the lists of opponent words and distractor words before providing a clue. I will also aim to create a clue that is more specific and less likely to be related to an opponent's word. For example, I could try to find a word that is related to a specific subset of my team's words, rather than trying to find a word that is related to multiple words. } \\
\\ 
\texttt{Additionally, I will pay closer attention to the format of my response, making sure to start with 'CLUE: ', followed by a new line, and then 'TARGETS: ' followed by a comma{-}separated list of related target words.} \\
\\ 
\texttt{To achieve this, I will take the following specific actions:} \\
\\ 
\texttt{1. Review the lists of opponent words and distractor words to identify potential pitfalls.} \\
\texttt{2. Identify a specific subset of my team's words that I want to relate to the clue.} \\
\texttt{3. Find a word that is semantically related to this subset, but not to any opponent's word or distractor word.} \\
\texttt{4. Ensure that my response follows the correct format, starting with 'CLUE: ', followed by a new line, and then 'TARGETS: ' followed by a comma{-}separated list of related target words.} \\
\\ 
\texttt{By following these steps, I hope to create a more effective clue that will help my teammate guess all of our team's words before the opposing team has uncovered all of theirs.} \\
\\ 
\\ 
\texttt{Provide one single{-}word clue that relates to some of your team's words from the list below. You can choose to how many and to which words your clue relates to.} \\
\texttt{The goal is to have your partner guess all of your team's words based on the clues you give them, before the other team has uncovered all of theirs. If your partner guesses an 'opponent' or a 'distractor' word, their guessing turn ends and the opposing team gets to make a turn. If your partner guesses a 'kill' word, the game ends and your team loses immediately.} \\
\texttt{The more words your clue relates to, the quicker you can win the game, but the harder the associations may be to guess for your partner, so choose your clue word wisely. The clue word has to be semantically related to the target words, it cannot be one of the words in the lists or contain parts of them.} \\
\\ 
\texttt{Always give your single{-}word clue and your comma{-}separated list of related target words in the following format and make your answers as short as possible, never include any other text than is required in this form:} \\
\\ 
\texttt{CLUE: <WORD>} \\
\texttt{TARGETS: <TARGETS>} \\
\\ 
\texttt{Your team words are: dilution, reinstatement, gracefulness, nationalization.} \\
\texttt{Your opponent's team words are: .} \\
\texttt{Distractor words are: stigma, ridiculousness, optimism, inefficiency, unsettlement, adorability.} \\
\texttt{The kill words are: coercion.} \\
            }
        }
    }
    & \\ \\

    \theutterance \stepcounter{utterance}  
    & \multicolumn{4}{p{0.6\linewidth}}{
        \cellcolor[rgb]{0.9,0.9,0.9}{
            \makecell[{{p{\linewidth}}}]{
                \texttt{\tiny{[P1$\rangle$GM]}}
                \texttt{CLUE: elegance} \\
\texttt{TARGETS: gracefulness, reinstatement, dilution} \\
            }
        }
    }
    & & \\ \\

    \theutterance \stepcounter{utterance}  
    & & & \multicolumn{2}{p{0.3\linewidth}}{
        \cellcolor[rgb]{0.9,0.9,0.9}{
            \makecell[{{p{\linewidth}}}]{
                \texttt{\tiny{[GM$|$GM]}}
                \texttt{elegance} \\
            }
        }
    }
    & & \\ \\

    \theutterance \stepcounter{utterance}  
    & & & \multicolumn{2}{p{0.3\linewidth}}{
        \cellcolor[rgb]{0.9,0.9,0.9}{
            \makecell[{{p{\linewidth}}}]{
                \texttt{\tiny{[GM$|$GM]}}
                ['gracefulness', 'reinstatement', 'dilution']
            }
        }
    }
    & & \\ \\

    \theutterance \stepcounter{utterance}  
    & & & \multicolumn{2}{p{0.3\linewidth}}{
        \cellcolor[rgb]{0.9,0.9,0.9}{
            \makecell[{{p{\linewidth}}}]{
                \texttt{\tiny{[GM$|$GM]}}
                {'word': 'gracefulness', 'assignment': 'team'}
            }
        }
    }
    & & \\ \\

    \theutterance \stepcounter{utterance}  
    & & & \multicolumn{2}{p{0.3\linewidth}}{
        \cellcolor[rgb]{0.9,0.9,0.9}{
            \makecell[{{p{\linewidth}}}]{
                \texttt{\tiny{[GM$|$GM]}}
                {'word': 'reinstatement', 'assignment': 'team'}
            }
        }
    }
    & & \\ \\

    \theutterance \stepcounter{utterance}  
    & & & \multicolumn{2}{p{0.3\linewidth}}{
        \cellcolor[rgb]{0.9,0.9,0.9}{
            \makecell[{{p{\linewidth}}}]{
                \texttt{\tiny{[GM$|$GM]}}
                {'word': 'dilution', 'assignment': 'team'}
            }
        }
    }
    & & \\ \\

    \theutterance \stepcounter{utterance}  
    & & & \multicolumn{2}{p{0.3\linewidth}}{
        \cellcolor[rgb]{0.9,0.9,0.9}{
            \makecell[{{p{\linewidth}}}]{
                \texttt{\tiny{[GM$|$GM]}}
                \texttt{opponent has won} \\
            }
        }
    }
    & & \\ \\

    \theutterance \stepcounter{utterance}  
    & \multicolumn{4}{p{0.6\linewidth}}{
        \cellcolor[rgb]{0.95,0.95,0.95}{
            \makecell[{{p{\linewidth}}}]{
                \texttt{\tiny{[P2$\rightarrow$P1]}}
                \texttt{Post{-}episode Reflexion (not sent to GM):} \\
\texttt{New plan: } \\
\\ 
\texttt{In my previous attempts, I provided a single{-}word clue that was semantically related to some of my team's words, but I failed to consider the potential for my clue to be related to an opponent's word, which led to my teammate's turn ending prematurely. } \\
\\ 
\texttt{Next time, I will make sure to carefully review the lists of opponent words and distractor words before providing a clue. I will also aim to create a clue that is more specific and less likely to be related to an opponent's word. For example, I could try to find a word that is related to a specific subset of my team's words, rather than trying to find a word that is related to multiple words. } \\
\\ 
\texttt{Additionally, I will pay closer attention to the format of my response, making sure to start with 'CLUE: ', followed by a new line, and then 'TARGETS: ' followed by a comma{-}separated list of related target words.} \\
\\ 
\texttt{To achieve this, I will take the following specific actions:} \\
\\ 
\texttt{1. Review the lists of opponent words and distractor words to identify potential pitfalls.} \\
\texttt{2. Identify a specific subset of my team's words that I want to relate to the clue.} \\
\texttt{3. Find a word that is semantically related to this subset, but not to any opponent's word or distractor word.} \\
\texttt{4. Ensure that my response follows the correct format, starting with 'CLUE: ', followed by a new line, and then 'TARGETS: ' followed by a comma{-}separated list of related target words.} \\
\\ 
\texttt{By following these steps, I hope to create a more effective clue that will help my teammate guess all of our team's words before the opposing team has uncovered all of theirs.} \\
            }
        }
    }
    & & \\ \\

\end{supertabular}
}

\end{document}
